\section{Training Configuration}

All methods are trained under a controlled adversarial training pipeline. Unless
otherwise stated, the optimizer, model architecture, dataset split, preprocessing
pipeline, and evaluation protocol are kept fixed across methods. This ensures that
performance differences mainly come from the training objective rather than from
implementation or evaluation differences.

To make method differences attributable to the training objective, Table~\ref{tab:ch5_training_config}
fixes the shared training configuration for the main experiments. All main experiments use ResNet-18 on CIFAR-10 with stochastic gradient descent (SGD),
momentum, weight decay, mixed-precision training, and a multi-step learning-rate
schedule. The main methods are trained for 50 epochs. During training, each minibatch
contains a balanced mixture of source attacks. In the two-source setting, approximately
half of the minibatch is assigned to PGD-\(\ell_\infty\) and the other half to
DDN-\(\ell_2\).

\begin{table}[H]
\centering
\caption{General training configuration used for the main experiments.}
\label{tab:ch5_training_config}
\footnotesize
\setlength{\tabcolsep}{4pt}
\renewcommand{\arraystretch}{1.08}
\begin{tabularx}{\textwidth}{@{}
>{\raggedright\arraybackslash}p{0.25\textwidth}
>{\raggedright\arraybackslash}p{0.22\textwidth}
>{\raggedright\arraybackslash}X
@{}}
\toprule
Component & Setting & Description \\
\midrule
Dataset & CIFAR-10 & 40k train / 10k validation / 10k test. \\
Backbone & ResNet-18 & CIFAR-10 adapted residual network. \\
Training epochs & 50 & Main training budget. \\
Batch size & 128 & Used for all main methods. \\
Optimizer & SGD & Momentum-based stochastic gradient descent. \\
Momentum & 0.9 & Standard SGD momentum. \\
Weight decay & $5 \times 10^{-4}$ & Weight regularization. \\
Initial learning rate & 0.1 & Decayed during training. \\
Learning-rate schedule & Multi-step & Same schedule for all compared methods. \\
Mixed precision & Enabled & Reduces memory usage and training time. \\
Gradient clipping & Enabled & Stabilizes adversarial training. \\
Checkpoint selection & Validation robust average & Test set is used only after selection. \\
\bottomrule
\end{tabularx}
\end{table}

Checkpoint selection is based on validation robustness rather than test performance.
After training, the selected checkpoint is evaluated once on the test set. This protocol
is applied consistently to the uniform baseline and all AttackDRO variants.

The method overview in Table~\ref{tab:ch5_method_training_summary} separates each
training objective by its grouping or weighting strategy. All multi-attack methods use
the same source attacks and mixed-batch adversarial example generation; they differ only
in how adversarial examples are weighted or grouped during optimization.

\begin{table}[H]
\centering
\caption{Training methods compared in the CIFAR-10 and ResNet-18 experiments.}
\label{tab:ch5_method_training_summary}
\footnotesize
\setlength{\tabcolsep}{4pt}
\renewcommand{\arraystretch}{1.08}
\begin{tabularx}{\textwidth}{@{}
>{\raggedright\arraybackslash}p{0.24\textwidth}
>{\raggedright\arraybackslash}p{0.34\textwidth}
>{\raggedright\arraybackslash}X
@{}}
\toprule
Method & Grouping / weighting strategy & Experimental role \\
\midrule
PGD-AT
& One source attack only
& Specialist $\ell_\infty$ baseline. \\

DDN-AT
& One source attack only
& Specialist $\ell_2$ baseline. \\

Multi-AT
& Uniform sample weighting over mixed PGD-$\ell_\infty$ and DDN-$\ell_2$ batches
& Main strong baseline. \\

AttackDRO
& Group DRO over attack identity
& Coarse attack-level DRO baseline. \\

Cluster-DRO
& Group DRO over discovered clusters
& Intermediate cluster-level grouping baseline. \\

AttackDRO++
& Uniform-anchored Cluster-DRO with gradient fingerprints
& Proposed final method. \\
\bottomrule
\end{tabularx}
\end{table}

For the proposed method, cluster assignments are refreshed periodically using a feature
bank. The default refresh interval is two epochs. The cluster weight vector
\(\mathbf{q}\) is updated only after a warmup period, so that early noisy representations
do not dominate adaptive reweighting.
